% Software Overview
% MIT License
%
% Copyright (c) 2026 HumberASV
% Copyright (c) 2026 Carson Fujita
%
% Permission is hereby granted, free of charge, to any person obtaining a copy
% of this software and associated documentation files (the "Software"), to deal
% in the Software without restriction, including without limitation the rights
% to use, copy, modify, merge, publish, distribute, sublicense, and/or sell
% copies of the Software, and to permit persons to whom the Software is
% furnished to do so, subject to the following conditions:
%
% The above copyright notice and this permission notice shall be included in all
% copies or substantial portions of the Software.
%
% THE SOFTWARE IS PROVIDED "AS IS", WITHOUT WARRANTY OF ANY KIND, EXPRESS OR
% IMPLIED, INCLUDING BUT NOT LIMITED TO THE WARRANTIES OF MERCHANTABILITY,
% FITNESS FOR A PARTICULAR PURPOSE AND NONINFRINGEMENT. IN NO EVENT SHALL THE
% AUTHORS OR COPYRIGHT HOLDERS BE LIABLE FOR ANY CLAIM, DAMAGES OR OTHER
% LIABILITY, WHETHER IN AN ACTION OF CONTRACT, TORT OR OTHERWISE, ARISING FROM,
% OUT OF OR IN CONNECTION WITH THE SOFTWARE OR THE USE OR OTHER DEALINGS IN THE
% SOFTWARE.

\subsection{Loon Env}
\label{sec:loon-env}

Loon-Env is the development environment for the Loon-E ASV project.

Clone it with submodules from GitHub:

\begin{verbatim}
    git clone --recurse-submodules https://github.com/HumberASV/Loon-Env.git
\end{verbatim}


The repository provides several scripts:

\begin{enumerate}
 \item \texttt{LoonEnv}: installation and update script.
 \item \texttt{LoonE}: wrapper command for building, starting, and entering containers.
 \item \texttt{LoonTools}: diagnostics and build helper commands.
 \item \texttt{configMTU}: script for MTU configuration related to ZEDX networking.
\end{enumerate}

\subsubsection{LoonEnv Command}

\texttt{LoonEnv} installs dependencies, sets environment variables, installs tools/assets, and prepares ROS 2 workspaces \autocite{LOONENV}.

Common usage patterns from the project workflow are:

\begin{verbatim}
    ./LoonEnv -v -s
\end{verbatim}

Installs or updates scripts and smaller assets (for example: \texttt{LoonTools}, \texttt{LoonE}, and config files).

\begin{verbatim}
    ./LoonEnv -v -f docker -s
\end{verbatim}

Restricts updates to Docker-related components.

\begin{verbatim}
    ./LoonEnv -v
\end{verbatim}

Runs the full installer, including MTU configuration and workspace setup.

\begin{verbatim}
    ./LoonEnv --dry-run
\end{verbatim}

Previews commands without applying changes.

\subsubsection{LoonE Command}

\texttt{LoonE} is the operator-facing command installed by Loon-Env for launching and interacting with the Docker-based Loon-E environment \autocite{LOONENV}.

The setup concepts from Section \ref{sec:docker-setup} and Section \ref{sec:env-vars-setup} are abstracted behind \texttt{LoonE} to provide a simpler interface for developers.

To view the available options:

\begin{verbatim}
    LoonE --help
\end{verbatim}

\subsubsection{What It Does}
\warningbox{LoonE wraps Docker and Docker Compose operations. Ensure Docker is installed and available in your terminal before using LoonE.}

\begin{itemize}
 \item Builds containers using hardware or virtual ZEDX profiles.
 \item Starts or enters specific services (for example, \texttt{zed} and \texttt{loone}).
 \item Uses installed assets and environment settings configured by \texttt{LoonEnv}.
\end{itemize}

Typical examples are:

\begin{verbatim}
    LoonE -b --virtual zed
    LoonE -b --hardware zed
    LoonE -b loone
    LoonE -s loone
    LoonE -e loone
\end{verbatim}

\subsubsection{Building with LoonE}

Use \texttt{-b} to build services:

\begin{verbatim}
    LoonE -b loone
\end{verbatim}

For ZED builds, select the profile explicitly:

\begin{verbatim}
    LoonE -b --virtual zed
    LoonE -b --hardware zed
\end{verbatim}

\subsubsection{Running with LoonE}

Start services with \texttt{-s}:

\begin{verbatim}
    LoonE -s loone
\end{verbatim}

Depending on local setup, developers may start one or more services as needed.

\subsubsection{Entering Containers}

Enter a running service container with \texttt{-e}:

\begin{verbatim}
    LoonE -e loone
\end{verbatim}

You can also use Docker directly when needed:

\begin{verbatim}
    docker ps
    docker exec -it loone bash
\end{verbatim}

\warningbox{The ZED container is intended for normal development as the non-root user. Use \texttt{sudo} only for maintenance tasks such as package refresh or dependency installation.}

\subsubsection{LoonTools Command}

\texttt{LoonTools} provides diagnostics and helper workflows.

Common options include:

\begin{verbatim}
    LoonTools --doctor
    LoonTools --colcon
    LoonTools --dry-run
    LoonTools --help
\end{verbatim}

\notebox{When \texttt{--colcon} is run in the ZED container, ROS setup is sourced in the same shell and the build runs as the normal container user.}

\subsubsection{configMTU Command}

\texttt{configMTU} configures MTU settings for network interfaces used by LoonE containers. This is important for reliable host-container communication with high-throughput ZEDX camera data.

\subsubsection{Environment Variables Installed by LoonEnv}

The LoonEnv setup configures standard variables such as:

\begin{itemize}
 \item \texttt{LOONE\_FILES\_DIR}: installation location for Loon-Env files (default: \texttt{/opt/loon\_env}).
 \item \texttt{LOONE\_WORKSPACE\_DIR}: ROS 2 workspace root (default: \texttt{\$HOME/LoonE\_ws}).
 \item \texttt{LOONE\_SETUP\_SCRIPT}: setup script path (default: \texttt{\$HOME/.loon\_env\_setup.bash}).
 \item \texttt{ROS\_DISTRO}: ROS distribution (default: \texttt{humble}).
 \item \texttt{RMW\_IMPLEMENTATION}: DDS implementation (default: \texttt{rmw\_cyclonedds\_cpp}).
\end{itemize}


\notebox{The generated ROS workspaces are mounted read-write, so normal development artifacts (for example, \texttt{colcon} build output) persist in the workspace.}
